\chapter{Full Measurement Results}
\label{app:results}

\Cref{tab:llr40-matrix} lists every cell of the loop-transformation matrix as the minimum of its
30 timed runs, in milliseconds, measured as described in \Cref{sec:setup}. The fastest valid
version of each kernel is set in bold.

\begin{table}[p]
  \centering
  \scriptsize
  \setlength{\tabcolsep}{4.5pt}
  % Generated by plotting/thesis/llr40_matrix.py -- do not edit.
\begin{tabular}{@{}lrrrrrr@{}}
  \toprule
  Kernel & C & Hand-written C & C++ & Fortran & Numba & Agent \\
  \midrule
  \multicolumn{7}{@{}l}{\emph{Interchange}} \\
  \quad\texttt{tsvc\_2\_s1232}$^\S$ & 39.1 & 56.2 & 50.1 & 55.6 & 83.5 & \textbf{13.1} \\
  \quad\texttt{tsvc\_2\_s2233} & 430.4 & 422.6 & 434.6 & \textbf{422.0} & 453.9 & --- \\
  \quad\texttt{tsvc\_2\_s231}$^\S$ & 723.3 & 547.5 & 694.0 & 706.7 & 594.7 & \textbf{69.9} \\
  \quad\texttt{tsvc\_2\_s233} & 681.6 & 690.9 & 681.9 & 657.1 & 673.5 & \textbf{66.7} \\
  \quad\texttt{tsvc\_2\_s235}$^\S$ & 789.2 & 799.2 & 545.0 & 539.4 & 460.2 & \textbf{58.0} \\
  \multicolumn{7}{@{}l}{\emph{Distribution}} \\
  \quad\texttt{tsvc\_2\_s2275}$^\S$ & 668.7 & 426.5 & 606.9 & 684.1 & 646.0 & \textbf{75.1} \\
  \multicolumn{7}{@{}l}{\emph{Fusion}} \\
  \quad\texttt{fuse\_diamond} & 168.8 & \textbf{29.8} & 169.3 & 167.8 & 186.4 & 30.9 \\
  \quad\texttt{fuse\_move\_ifs} & 68.2 & 67.9 & 67.7 & 67.5 & 65.3 & \textbf{51.7} \\
  \quad\texttt{fuse\_stencil\_through\_transient} & 112.3 & \textbf{47.7} & 113.0 & 114.4 & 186.3 & 47.8 \\
  \multicolumn{7}{@{}l}{\emph{Control Flow}} \\
  \quad\texttt{ext\_break\_capture} & 67.1 & 66.9 & 66.4 & 66.9 & 70.9 & \textbf{29.7} \\
  \quad\texttt{tsvc\_2\_s2710} & 220.6 & 219.7 & 216.7 & 225.1 & 93.4 & \textbf{84.5} \\
  \quad\texttt{tsvc\_2\_s275}$^\S$ & 704.0 & 720.8 & 719.9 & 545.7 & 476.2 & \textbf{163.9} \\
  \quad\texttt{tsvc\_2\_vpvts} & 88.3 & 87.3 & 89.5 & 89.6 & \textbf{87.1} & 87.1 \\
  \quad\texttt{tsvc\_2\_vtvtv} & 96.4 & 94.9 & 96.3 & 95.8 & 94.8 & \textbf{94.8} \\
  \multicolumn{7}{@{}l}{\emph{Reduction}} \\
  \quad\texttt{argmax\_with\_index} & 71.6 & \textbf{71.5} & 71.5 & 71.7 & 178.6 & 71.7 \\
  \quad\texttt{quasi\_affine\_reduce\_odd} & 63.8 & 66.3 & 63.6 & 63.6 & 67.4 & \textbf{63.5} \\
  \quad\texttt{scan\_affine\_decay} & \textbf{118.2} & --- & 118.3 & 148.0 & 147.5 & 235.3 \\
  \quad\texttt{tsvc\_2\_s311} & 135.5 & 135.7 & 135.5 & 135.4 & 135.4 & \textbf{76.2} \\
  \quad\texttt{tsvc\_2\_s3110} & \textbf{120.2} & 120.8 & 120.4 & 121.5 & 315.3 & 121.9 \\
  \quad\texttt{tsvc\_2\_s3111} & 126.5 & 126.3 & \textbf{124.9} & 728.6 & 314.0 & 125.8 \\
  \quad\texttt{tsvc\_2\_s3112} & 79.8 & 79.8 & 79.9 & 79.8 & \textbf{79.7} & 79.9 \\
  \quad\texttt{tsvc\_2\_s316} & 275.8 & 276.0 & 80.0 & 80.2 & 123.1 & \textbf{69.3} \\
  \quad\texttt{tsvc\_2\_s318} & 80.2 & 80.2 & \textbf{79.7} & 80.0 & 290.7 & 252.6 \\
  \quad\texttt{tsvc\_2\_s319} & 75.0 & 74.9 & 74.5 & 74.8 & 76.6 & \textbf{74.1} \\
  \multicolumn{7}{@{}l}{\emph{Recurrence}} \\
  \quad\texttt{tsvc\_2\_s323} & 83.0 & 82.8 & 88.2 & \textbf{80.1} & 82.0 & 82.7 \\
  \multicolumn{7}{@{}l}{\emph{Scalar Exp.}} \\
  \quad\texttt{tsvc\_2\_s252} & 82.3 & 85.9 & 87.7 & 85.6 & \textbf{80.7} & 83.9 \\
  \quad\texttt{tsvc\_2\_s255} & 77.6 & 77.5 & 77.3 & 77.7 & 80.6 & \textbf{66.1} \\
  \multicolumn{7}{@{}l}{\emph{Linear Dep.}} \\
  \quad\texttt{tsvc\_2\_s115} & 83.1$^\ddagger$ & 84.0$^\ddagger$ & 83.3$^\ddagger$ & 84.3$^\ddagger$ & \textbf{141.0} & 86.8$^\ddagger$ \\
  \quad\texttt{tsvc\_2\_s119} & 70.1 & 70.5 & 70.4 & 70.0 & 69.1 & \textbf{56.9} \\
  \multicolumn{7}{@{}l}{\emph{Dep. Dist.}} \\
  \quad\texttt{ext\_war\_unit} & 61.6 & 61.3 & 61.5 & \textbf{60.4} & 61.5 & 61.5 \\
  \quad\texttt{versioned\_distance\_update} & 302.9 & --- & 300.6 & 330.8 & 325.2 & \textbf{123.1} \\
  \multicolumn{7}{@{}l}{\emph{Node Split.}} \\
  \quad\texttt{tsvc\_2\_s1244} & 53.9 & \textbf{53.1} & 53.6 & 53.6 & 53.5 & 58.0 \\
  \multicolumn{7}{@{}l}{\emph{Interproc.}} \\
  \quad\texttt{tsvc\_2\_s152} & 63.2 & \textbf{51.5} & 62.5 & 62.6 & 63.1 & 51.6 \\
  \multicolumn{7}{@{}l}{\emph{Indirect}} \\
  \quad\texttt{scatter\_accum\_dup} & \textbf{351.2} & --- & 352.4 & 371.4 & 409.8 & 938.4 \\
  \quad\texttt{segment\_reduce\_ragged} & 39.0 & --- & 38.9 & 39.8 & 56.1 & \textbf{38.0} \\
  \quad\texttt{tsvc\_2\_s4112} & 465.5 & 463.7 & 463.0 & 486.2 & 468.9 & \textbf{449.7} \\
  \quad\texttt{tsvc\_2\_vag} & 447.8 & 439.4 & 449.6 & 464.6 & \textbf{438.2} & 482.5 \\
  \multicolumn{7}{@{}l}{\emph{Packing}} \\
  \quad\texttt{compact\_threshold\_pack} & 110.9 & --- & 111.4 & 111.1 & \textbf{106.2} & 143.7 \\
  \multicolumn{7}{@{}l}{\emph{Wavefront}} \\
  \quad\texttt{wf\_diff\_skew}$^\dagger$ & \textbf{59.9} & 60.4 & 60.1 & 60.1 & 61.5 & 66.3 \\
  \quad\texttt{wf\_triangular}$^\dagger$ & 99.5 & 100.5 & 99.7 & \textbf{99.5} & 382.8 & 153.4 \\
  \bottomrule
\end{tabular}

  \caption{The loop-transformation matrix: minimum of 30 runs in milliseconds, preset M,
    \texttt{float64}, one core of a GH200 node. Bold: fastest valid version. ---: version does
    not exist. $^\ddagger$: output does not match the NumPy oracle; the time is shown but not
    used. $^\dagger$: passes validation only vacuously at this size
    (\Cref{sec:language-comparison}). $^\S$: GCC versions not reproducible between two
    measurements; their differences are memory-bandwidth noise (\Cref{sec:setup}).}
  \label{tab:llr40-matrix}
\end{table}

\chapter{Extracted Kernel Catalogue}
\label{app:extracted}
% Per kernel: upstream file and commit, licence, presets, graded outputs, PR.

\chapter{Skill Texts}
\label{app:skills}
% Kernel-extraction and solver-preconditioning skills, or pointers to the repository.

\chapter{Reproducibility}
\label{app:reproduce}
% Pinned commits, environment record, how to rerun each measurement.
% LLR-40 figure and tables: python plotting/run.py plotting/thesis/llr40_matrix.py
